\begin{textsample}{POS Dim 6 – se – Score 10.00 – se/se\_em\_29.txt}  \label{ex:f6_pos_042}
5.6 Oferta Desde 2007, a Educação Profissional e Tecnológica ( EPT ) é uma modalidade educacional oferecida na esfera estadual.
Esta modalidade está distribuída regionalmente e envolve uma Rede de Centros Estaduais de Educação Profissional e Instituições educacionais compartilhadas, nas Diretorias Regionais de Ensino, abrangendo quase todas as regiões geoeconômicas do estado, segundo o mapa a seguir.
A oferta envolve uma Rede de 09 Centros Estaduais de Educação Profissional e cinco Instituições educacionais compartilhadas, oferecendo até 2020, 30 cursos técnicos, nas mais diversas formas ( integrado, concomitante e subsequente ) e 22 cursos FIC de acordo com \textbf{portfólio} em anexo.
Este \textbf{Portfólio} tem como propósito apresentar um leque de ofertas não exaustivo como mecanismo de direcionar toda comunidade escolar. 5.7 Formação Básica para o Trabalho e Eixos \textbf{Estruturantes} A Educação Profissional tem sido foco de estudos por todos os segmentos da sociedade organizada, face às mudanças que vêm ocorrendo no cenário futuro de trabalho com o surgimento de novas profissões e necessidade de novos talentos.
O sentido de educação “ consiste naquilo que ela carrega de projeção do futuro ”. ( Sobrinho, 1995 ).
Para TOFLLER ( 1992:25 ) “ toda educação parte de uma imagem do futuro.
Se a ideia do futuro que uma sociedade tem é toscamente inadequada, seu sistema educativo atraiçoará sua juventude ”.
A organização curricular dos cursos técnicos é pautada, em primeiro lugar, pela definição do \textbf{PERFIL} PROFISSIONAL de conclusão, tornando a concepção curricular, um meio pedagógico indispensável, na organização dos currículos de educação profissional.
O artigo 17 da Resolução 6/2012 reforça : > “ O \textbf{perfil} profissional de conclusão do curso, é definido pela explicitação dos conhecimentos, saberes e competências profissionais e pessoais, tanto aquelas que caracterizam a preparação básica para o trabalho, quanto as comuns para o respectivo eixo tecnológico, bem como as específicas de cada habilitação profissional e das etapas de qualificação e de especialização profissional técnica que compõem o correspondente itinerário \textbf{formativo} ”.
As DNCEM propõem que os Itinerários \textbf{Formativos} trabalhem unidades curriculares e atividades integradoras, organizadas por meio da preparação básica para o trabalho e quatro eixos \textbf{estruturantes}, inerentes à ocupação/curso escolhido.
Em conformidade com os eixos \textbf{estruturantes}, nos cursos técnicos, foi incluído o componente curricular denominado Projeto de Aprendizagem Interdisciplinar - PAI, que integra conhecimentos voltados para rotinas do eixo tecnológico, o mundo do trabalho e a prática da ocupação estudada.
São Eixos \textbf{Estruturantes} : - Investigação Científica – integra o componente curricular PROJETO DE APRENDIZAGEM INTERDISCIPLINAR - PAI I, encaminhando o aluno para a realização de procedimentos de investigação voltados à compreensão e enfrentamento de situações, com proposição de intervenções que considerem o desenvolvimento local e a melhoria da qualidade de vida da comunidade.
Define -se o \textbf{problema} de pesquisa ; os objetivos ; através da estrutura de trabalhos científicos ; - Processos Criativos - que integra o componente curricular Projeto de aprendizagem interdisciplinar, PAI II.
Procura -se expandir a capacidade dos estudantes de idealizar e realizar projetos criativos ; a partir dos conhecimentos adquiridos no decorrer do curso ; - \textbf{Mediação} e Intervenção Sociocultural - No componente curricular Projeto de aprendizagem interdisciplinar, PAI III, os estudantes são envolvidos em campos de atuação da vida pública, por meio de projeto que os levem a promover transformações positivas na comunidade, ampliando habilidades relacionadas à convivência e atuação sociocultural ; utilizando esses conhecimentos e habilidades para mediar conflitos, promover entendimentos e propor soluções para questões e \textbf{problemas} socioculturais e \textbf{ambientais} identificados nas comunidades. - \textbf{Empreendedorismo} - Com foco no desenvolvimento de processos e produtos com o uso de \textbf{tecnologias} variadas.
O estudante visualiza as características do comportamento empreendedor e sua importância para o desenvolvimento pessoal e profissional, aplica modelos mentais e técnicas de desenvolvimento do \textbf{perfil} empreendedor, identifica oportunidades para empreender do mercado, articula competências gerais do curso para construção do Projeto.
Os eixos \textbf{estruturantes} são complementares e é importante que os Itinerários \textbf{Formativos} incorporem e integrem todos eles, a fim de garantir que os estudantes experimentem diferentes situações de aprendizagem e desenvolvam um conjunto diversificado de habilidades relevantes para sua formação integral.
A Educação Profissional e tecnológica, tem o poder de contribuir com o aumento da capacidade de (re)inserção social, laboral e política dos seus formandos ; com a extensão de ofertas que contribuam à formação integral dos coletivos que procuram a instituição educacional pública de EPT para que esses sujeitos possam atuar, de forma competente e ética, como agentes de mudanças orientadas à satisfação das necessidades coletivas, notadamente as das classes trabalhadoras ( MOURA, 2000 ; FREIRE, 1986 ; 2000a e b ; 2001 ).
Assim, os estudantes, no decorrer de seu Ensino Médio, deverão realizar pelo menos um itinerário \textbf{formativo} completo, passando, necessariamente por todos os eixos.

% matched lemmas: ambiental, empreendedorismo, estruturante, formativo, mediação, perfil, portfólio, problema, tecnologia
\end{textsample}
